% vim: ai sw=2

\documentclass[handout,xcolor=dvipsnames]{beamer}

\let\Oldemph\emph
\renewcommand{\emph}[1]{\textcolor{blue}{\Oldemph{#1}}}
\newcommand{\heading}[1]{\textbf{\textcolor{BurntOrange}{#1}}}

% Copyright 2004 by Till Tantau <tantau@users.sourceforge.net>.
%
% In principle, this file can be redistributed and/or modified under
% the terms of the GNU Public License, version 2.
%
% However, this file is supposed to be a template to be modified
% for your own needs. For this reason, if you use this file as a
% template and not specifically distribute it as part of a another
% package/program, I grant the extra permission to freely copy and
% modify this file as you see fit and even to delete this copyright
% notice. 


\mode<presentation>
{
%  \usetheme{Warsaw}
\usetheme{Madrid}

  \setbeamercovered{transparent}
}

\usepackage[utf8]{inputenc}
\usepackage{times}
\usepackage[T1]{fontenc}

\usepackage{graphicx}
\usepackage[final]{listings}
\usepackage{multirow}

\lstset{tabsize=2,numbers=none,showstringspaces=false,breaklines=true,
	breakatwhitespace=true,escapeinside={(*@}{@*)},
	basicstyle=\scriptsize,keywordstyle=\bfseries\color{OliveGreen},
	commentstyle=\color{blue}}

\title[Kernel Bug-finding] % (optional, use only with long paper titles)
{Automatic Bug-finding Techniques for Linux Kernel}

\subtitle
{Obhajoba tezí disertační práce}

\author[J. Slabý]{\emph{Jiří~Slabý}}

\institute[FI, MU] % (optional, but mostly needed)
{Fakulta Informatiky, Masarykova Univerzita}

\date[24. 5. 2010, Brno]{Brno, 24. května 2010}

\subject{Theoretical Computer Science}

\begin{document}

\begin{frame}
  \titlepage
\end{frame}

%%%%%%%%%%%%%%%%%%%%%%%

\begin{frame}
  \frametitle{Struktura obhajoby}
  \begin{itemize}
  \item Záměr práce
  \item Stav práce
  \item Vyjádření k posudkům
  \end{itemize}
\end{frame}

\begin{frame}
  \frametitle{Záměr práce -- Co a proč?}
  \begin{itemize}
  \item Automatická kontrola kódu
  \begin{itemize}
  \item Lidé se nechtějí učit složité nástroje
  \item Vstup: \emph{kód} a popř. definice kontrol (jinak možná autodetekce)
  \item Výstup: \emph{seznam chyb} (co nejméně falešných)
  \end{itemize}\pause
  \item Podpora jádra
  \begin{itemize}
  \item Chyby v jádře končí katastrofou
  \item Obsáhlý kód, rozmanitost
  \end{itemize}\pause
  \item Podpora Linuxového jádra
  \begin{itemize}
  \item Je rošířené, testované, přesto jsou v něm chyby -- výzva
  \item Ignoruje uváznutí, čeká že nenastane
  \end{itemize}
  \end{itemize}

\end{frame}

\begin{frame}
  \frametitle{Záměr práce -- Jak?}
  \begin{itemize}
  \item To ví Bůh
  \end{itemize}

\end{frame}

\begin{frame}
  \frametitle{Stav práce}
  \begin{itemize}
  \item Učící fáze: \textsc{Stanse}
  \begin{itemize}
  \item Statická analýza, počítá fix-point abstraktní interpretace
  \item Podpora jádra Linuxu, nalezeno několik desítek chyb
  \end{itemize}
  \item Nyní nový nástroj
  \begin{itemize}
  \item Založený na \textsc{Stanse}, \textsc{Sparse}, \textsc{Klee} a
  \textsc{LLVM}
  \end{itemize}
  \end{itemize}

\end{frame}

\begin{frame}
  \frametitle{Připomínky v posudcích}
  \begin{itemize}
  \item Špatné citace
  \item Přesný význam "Mixed analysis", Tabulka 2
  \item Kombinace SE a AI
  \item Referování chyb v kódu
  \item Žádné publikace
  \item Prezentace výsledků
  \end{itemize}

\end{frame}

\begin{frame}
  \frametitle{Závěr}
  \begin{itemize}
  \item Navržený směr disertační práce je jasný a splnitelný
  \item Nový nástroj se drží tématu tezí
  \item Termíny dodrženy
  \end{itemize}

\end{frame}

\begin{frame}
  \frametitle{Poděkování}
  \begin{center} \LARGE
  Děkuji za pozornost!
  \end{center}

\end{frame}

\end{document}

